%%%%%%%%%%%%%%%%%%%%%%%%%%%%%%%%%%%%%%%%%%%%%%%%%%%%%%%%%%%%%%%%%%%%%%
% Lớp 10 | Véctơ (chưa xét tọa độ) | Tổng và hiệu của hai véctơ
% Dạng 0H5?2-1: Tính toán, thu gọn hiệu các véctơ
% Số câu: 10  (sinh tự động bởi tach_cau.py; từ nay thầy sửa trực tiếp file này)
%%%%%%%%%%%%%%%%%%%%%%%%%%%%%%%%%%%%%%%%%%%%%%%%%%%%%%%%%%%%%%%%%%%%%%

%%%%%%%%%%%%%%%% TRẮC NGHIỆM 4 PHƯƠNG ÁN %%%%%%%%%%%%%%%%

% nguồn: MaHoaNH_TracNghiem_Lop10
\begin{ex}%[0H5N2-1]%[ma:5fc7be91]
	Trong mặt phẳng $ Oxy $, cho  hai vecto $ \vec{a} = (2;-4) $, $ \vec{b} = (-5;3) $. Tọa độ vecto $ \vec{u} = 2\vec{a} - \vec{b} $ là
	 \choice
		{$(-1;5)$} 
		{$(7;-7)$} 
		{\True  $(9;-11)$}
		{$(9;5)$} 
	\loigiai{ 
		Ta có $ \vec{u} = 2\vec{a} - \vec{b} = (4;-8) -  (-5;3) = (4+5; -8-3) = (9;-11).$\\
		Vậy tọa độ vecto $ \vec{u} $  là $ (9;-11) $.
	 }
\end{ex}

% nguồn: MaHoaNH_TracNghiem_Lop10
\begin{ex}%[0H5H2-1]%[ma:2d838667]
	Cho hình bình hành $ABCD$ tâm $O$. Khi đó $\overrightarrow{OA}+\overrightarrow{BO}$  bằng
	 \choice
		 {$\overrightarrow{OC}+\overrightarrow{OB}$}
		{$\overrightarrow{AB}$} 
		 {\True $\overrightarrow{CD}$}
		{$\overrightarrow{OC}+\overrightarrow{DO}$} 
	 \loigiai{
	\immini{ 
		Ta có $\overrightarrow{OA}+\overrightarrow{BO}=\overrightarrow{BO}+\overrightarrow{OA}=\overrightarrow{BA}$.\\ 
		Vì $ABCD$ là hình bình hành nên  $\overrightarrow{BA}=\overrightarrow{CD}$.
		}{\begin{tikzpicture}[line join=round,line cap=round,>=stealth,scale=.75,font=\footnotesize]
		 \foreach \x/\y/\n in {0/0/A,1/2/B,4/0/D} \coordinate (\n) at (\x,\y);
		\coordinate (C)  at ($(B)+(D)-(A)$);
		\coordinate (O)  at ($(B)!0.5!(D)$);
		\draw (A)--(B)--(C)--(D)--(A)--(C) 
			(B)--(D); 
		\foreach \t/\g  in {A/-135,C/45,B/135,D/-45,O/-90}{
			\draw[fill=black] (\t)  circle (1pt) node[shift={(\g:8pt)}]{$ \t $};
		} 
	\end{tikzpicture}
		} 
	} 
\end{ex}


%%%%%%%%%%%%%%%% ĐÚNG/SAI %%%%%%%%%%%%%%%%

% nguồn: MaHoaNH_DungSai_Lop10
\begin{ex}%[0H5N2-1]%[ma:392b58d7]
	Cho ba điểm bất kỳ $A$, $B$, $C$. Hãy  viết tính đúng sai của các câu sau.
	\choiceTF 
		{\True  $\vec{AB}=\vec{CB}-\vec{CA}$}
		{$\vec{BC}=\vec{AB}-\vec{AC}$} 
		 {$\vec{AC}=\vec{CB}-\vec{BA}$}
		{\True  $\vec{CA}=\vec{CB}+\vec{BA}$}
	\loigiai{ 
		\begin{itemize}
			\item  Hệ thức \fbox{$\vec{CB}- \vec{CA}=\vec{AB}$} luôn đúng với mọi điểm $A$, $B$, $C$.
			 \item Hệ thức $\vec{BC}=\vec{AB}- \vec{AC}=\vec{CB}$ chỉ đúng khi $B\equiv C$.
			\item Hệ thức $\vec{AC}=\vec{CB}- \vec{BA}\Leftrightarrow \vec{BA}+\vec{AC}=\vec{CB}\Leftrightarrow \vec{BC}=\vec{CB}$ chỉ đúng khi  $B\equiv C$.
			\item  Hệ thức $\vec{CA}=\vec{CB}+ \vec{BA}$ luôn đúng với mọi điểm $A$, $B$, $C$.
		\end{itemize}
	} 
\end{ex}

% nguồn: MaHoaNH_DungSai_Lop10
\begin{ex}%[0H5N2-1]%[ma:c8fb2e02]
	Gọi $O$ là tâm  của hình vuông $ABCD$. Véc-tơ nào trong các véc-tơ dưới đây bằng $\vec{CA}$
	 \choiceTF
		{$\vec{BC}+\vec{AB}$} 
		{$\vec{OA}+\vec{OC}$} 
		 {\True $\vec{BA}+\vec{DA}$}
		{\True $\vec{CD}-  \vec{BC}$}
	\loigiai{ 
		\begin{itemchoice}
			 \itemch $\vec{BC}+\vec{AB}=\vec{AB}+\vec{BC}=\vec{AC}$.
			\itemch $\vec{OA}+\vec{OC}=\vec{0}$. 
			\itemch $\vec{BA}+\vec{DA}=\vec{CB}+\vec{BA}=\vec{CA}$. 
			 \itemch $\vec{CD}- \vec{BC}=\vec{CD}+\vec{CB}=\vec{CA}$.
		\end{itemchoice}
	} 
\end{ex}

% nguồn: MaHoaNH_DungSai_Lop10
\begin{ex}%[0H5N2-1]%[ma:72285b31]
	Cho hình bình hành $ABCD$. Gọi $O$ là giao điểm của $AC$ và $BD$. Hãy  viết tính đúng sai của các câu sau.
	\choiceTF 
		{\True $\vec{OA}$ và $\vec{OC}$ là hai vec-tơ đối  nhau}
		{$\vec{OB}$ và  $\vec{DO}$ là hai vec-tơ đối nhau}
		{\True  $\vec{OA}+\vec{OC}=\vec{OB}+\vec{OD}$}
		 {$\vec{OA}+\vec{OB}+\vec{OC}+\vec{OD}=\vec{AD}+\vec{CD}$}
	\loigiai{ 
	\begin{itemchoice}
		\itemch $\vec{OA}$ và $\vec{OC}$ là hai vec-tơ đối nhau. 
		\itemch $\vec{OB}$  và $\vec{DO}$ là hai vec-tơ không đối nhau.
		\itemch  $\vec{OA}+\vec{OC}=\vec{OB}+\vec{OD}=\vec{0}$.
		\itemch $\vec{OA}+\vec{OB}+\vec{OC}+\vec{OD}\ne  \vec{AD}+\vec{CD}$
	\end{itemchoice}
	} 
\end{ex}

% nguồn: MaHoaNH_DungSai_Lop10
\begin{ex}%[0H5N2-1]%[ma:dc37e731]
	Với bốn điểm $A$, $B$, $C$, $D$ trong đó không có ba điểm nào thẳng hàng. Hãy viết tính đúng sai của  các câu sau.
	\choiceTF 
		{\True $ABCD$ là hình bình hành khi  $\vec{AB}=\vec{DC}$}
		{\True $ABCD$ là hình bình  hành khi $\vec{AB}+\vec{AD}=\vec{AC}$}
		{$ABCD$ là hình bình  hành khi $\vec{AC}=\vec{BD}$}
		{$ABCD$ là hình  bình hành khi $\vec{AD}=\vec{CB}$}
	 \loigiai{
		\begin{itemize}
			\item $\vec{AB}=\vec{DC}\Rightarrow  \heva{&AB\parallel DC\\&AB=DC.}$\\
			Vậy $ABCD$ là  hình bình hành.
		 \immini{
			\item $\vec{AD}=\vec{BC}\Rightarrow  \heva{&AD\parallel BC\\&AD=BC.}$\\
			 Vậy $ABCD$ là hình bình hành.
			 \item $\vec{AB}+\vec{AD}=\vec{AC}\Leftrightarrow \vec{AB}=\vec{AC}-\vec{AD}=\vec{DC}$.\\
			 Vậy $ABCD$ là hình bình hành.
			 }
			{ 
			\begin{tikzpicture}[scale=1,line join=round,line cap=round]
					 \tkzDefPoints{0/0/A,4/0/B,0.5/2/D}
					\tkzDefPointBy[translation  = from A to B](D) \tkzGetPoint{C}
					\tkzDrawPolygon  (A,B,C,D)
					\tkzLabelPoints[above](C,D) 
					\tkzLabelPoints[below](A,B) 
					\tkzDrawPoints (A,B,C,D) 
				\end{tikzpicture}
			 }
			 \item $\vec{AD}=\vec{CB}\Rightarrow \heva{&AD\parallel CB\\&AD=CB}$\\
			$\Rightarrow$ ta thu  được hình bình hành $ABDC$, \fbox{không phải là hình bình hành $ABCD$}.
		\end{itemize}
	 }
\end{ex}


%%%%%%%%%%%%%%%% TỰ LUẬN %%%%%%%%%%%%%%%%

% nguồn: MaHoaNH_TuLuan_Lop10
\begin{bt}%[0H5H2-1]%[ma:ee02117c]
	Cho hình bình hành $ABCD$ với $M$ và $N$ lần lượt là trung điểm của $BC$ và $AD$.  Tìm tổng của hai vectơ
	\begin{enumEX}{2}
		 \item $\overrightarrow{NC}$ và $\overrightarrow{MC}$.
		\item $\overrightarrow{AM}$ và  $\overrightarrow{CD}$.
	\end{enumEX}
	\loigiai{ 
		\begin{center}
			\begin{tikzpicture}[scale=1,font=\footnotesize,line join=round,line cap=round,>=stealth]
					 \def\a{4}
					 \path (0:0) coordinate (A)
						(0:\a) coordinate  (D)
						(65:.6*\a) coordinate  (B)
						($(D)-(A)+(B)$) coordinate  (C)
						($(C)!.5!(B)$) coordinate  (M)
						($(A)!.5!(D)$) coordinate (N); 
					 \draw (A)-- (B)-- (C)--(D)--cycle (A)-- (M)
						(C)-- (N); 
					 \foreach \x/\g in {A/180,B/150,C/45,D/0,M/90,N/-90}
					\fill[black] (\x) circle (1pt) 
						($(\g:3mm)+(\x)$) node  {$\x$};
				\end{tikzpicture}
		\end{center}
		\begin{enumerate}
			 \item Vì $\overrightarrow{MC}=\overrightarrow{AN}$ nên ta có $\overrightarrow{NC}+\overrightarrow{MC}=\overrightarrow{NC}+\overrightarrow{AN}=\overrightarrow{AN}+\overrightarrow{NC}=\overrightarrow{AC}$.
			\item  Vì $\overrightarrow{CD}=\overrightarrow{BA}$ nên ta có $\overrightarrow{AM}+\overrightarrow{CD}=\overrightarrow{AM}+\overrightarrow{BA}=\overrightarrow{BA}+\overrightarrow{AM}=\overrightarrow{BM}$.
		\end{enumerate}
	 }
\end{bt}

% nguồn: MaHoaNH_TuLuan_Lop10
\begin{bt}%[0H5H2-1]%[ma:bded88ed]
	Cho  tam giác $ABC$. Các điểm $M$, $N$ và $P$ lần lượt là trung điểm các cạnh $AB$, $AC$ và $BC$.
	\begin{enumerate}
		\item Tìm các hiệu sau $\overrightarrow{AM}-\overrightarrow{AN}$;  $\overrightarrow{MN}-\overrightarrow{NC}$ và $\overrightarrow{MN}-\overrightarrow{PN}$.
		\item Phân tích vectơ  $\overrightarrow{AM}$ theo hai vectơ $\overrightarrow{MN}$ và $\overrightarrow{MP}$.
	\end{enumerate}
	\loigiai{ 
		\begin{center}
			\begin{tikzpicture}[scale=1,font=\footnotesize,line join=round,line cap=round,>=stealth]
					\def\a{4} 
					\path  (0:0) coordinate (B)
						(0:\a)  coordinate (C)
						(65:.85*\a)  coordinate (A)
						($(A)!.5!(B)$) coordinate  (M)
						($(A)!.5!(C)$)  coordinate (N)
						($(B)!.5!(C)$) coordinate  (P);
					\draw (A)-- (B)--  (C)--cycle (M)-- (N)-- (P)--cycle;
					\foreach \x/\g in {A/90,B/150,C/45,M/145,N/45,P/-90} 
					\fill[black] (\x) circle (1pt) 
						($(\g:3mm)+(\x)$) node {$\x$}; 
				\end{tikzpicture}
		\end{center}
		\begin{enumerate}
			\item \begin{itemize}
			\item Theo  quy tắc ba điểm, thì $\overrightarrow{AM}-\overrightarrow{AN}=\overrightarrow{NM}$.
			\item Ta có $MP$ là đường trung bình của tam giác $ABC$ và $\overrightarrow{MP}$ cùng hướng với $\overrightarrow{NC}$ nên ta có $\overrightarrow{NC}=\overrightarrow{MP}$.  Do vậy $\overrightarrow{MN}-\overrightarrow{NC}=\overrightarrow{MN}-\overrightarrow{MP}=\overrightarrow{PN}$.
			\item Ta  có $-\overrightarrow{PN}=\overrightarrow{NP}$ nên $\overrightarrow{MN}-\overrightarrow{PN}=\overrightarrow{MN}+\overrightarrow{NP}=\overrightarrow{MP}$.
		\end{itemize}
		\item Ta có $\overrightarrow{AM}=\overrightarrow{NP}$ nên có phân  tích sau $\overrightarrow{AM}=\overrightarrow{NP}=\overrightarrow{MP}-\overrightarrow{MN}$.
	\end{enumerate}
} 
\end{bt}

% nguồn: MaHoaNH_TuLuan_Lop10
\begin{bt}%[0H5V2-1]%[ma:422c661e]
	Cho $n$ điểm $A_1$, $A_2$, $A_3$, $\dots$, $A_n$, xác định vectơ \[\overrightarrow{A_{n-1} A_n}+\overrightarrow{A_{n-2} A_{n-1}}+\overrightarrow{A_{n-3} A_{n-2}}+\cdots+\overrightarrow{A_2 A_3}+\overrightarrow{A_1 A_2}.\] 
	 \loigiai{
		Ta có \begin{eqnarray*}
		&&\overrightarrow{A_{n-1} A_n}+\overrightarrow{A_{n-2} A_{n-1}}+\overrightarrow{A_{n-3} A_{n-2}}+\cdots+\overrightarrow{A_2 A_3}+\overrightarrow{A_1  A_2}\\
		&=&\overrightarrow{A_1 A_2}+\overrightarrow{A_2 A_3}+\cdots+\overrightarrow{A_{n-3} A_{n-2}}+\overrightarrow{A_{n-2} A_{n-1}}+\overrightarrow{A_{n-1} A_n}. 
	\end{eqnarray*}
	Do đó $\overrightarrow{A_{n-1} A_n}+\overrightarrow{A_{n-2} A_{n-1}}+\overrightarrow{A_{n-3} A_{n-2}}+\cdots+\overrightarrow{A_2 A_3}+\overrightarrow{A_1 A_2}=\overrightarrow{A_1 A_n}$. 
 }
\end{bt}

% nguồn: MaHoaNH_TuLuan_Lop10
\begin{bt}%[0H5V2-1]%[ma:945fd082]
	 Cho tam giác $ABC$. Gọi $M$, $N$, $P$ lần lượt là trung điểm của các cạnh $BC$, $CA$, $AB$. Chứng minh rằng
	\begin{enumerate}
		 \item $\overrightarrow{BM}+\overrightarrow{CN}+\overrightarrow{AP}=\overrightarrow0$.
		\item  $\overrightarrow{OA}+\overrightarrow{OB}+\overrightarrow{OC}=\overrightarrow{OM}+\overrightarrow{ON}+\overrightarrow{OP}$, với $O$ là điểm bất kì.
	\end{enumerate}
	\loigiai{ 
		\begin{center}
			\begin{tikzpicture}[scale=1,font=\footnotesize,line join=round,line cap=round,>=stealth]
					 \def\a{4}
					\path (0:0) coordinate (B) 
						(0:\a)  coordinate (C)
						(65:.85*\a) coordinate (A) 
						($(A)!.5!(B)$) coordinate  (P)
						($(A)!.5!(C)$)  coordinate (N)
						 ($(B)!.5!(C)$) coordinate (M);
					\draw (A)-- (B)-- (C)--cycle (M)--  (N)-- (P)--cycle;
					\foreach \x/\g in  {A/90,B/150,C/45,M/-90,N/45,P/150}
					\fill[black]  (\x) circle (1pt)
						($(\g:3mm)+(\x)$) node  {$\x$};
				\end{tikzpicture}
		\end{center}
		\begin{enumerate}
			\item Ta có $PN$, $MN$ là đường trung bình của tam giác $ABC$ nên $\heva{&PN\parallel BM\\&MN\parallel BP.}$\\ 
			Suy ra tứ giác $BMNP$ là hình bình hành.\\ 
			Do đó, $\overrightarrow{BM}=\overrightarrow{PN}$ và $N$ là trung điểm của  $AC\Rightarrow \overrightarrow{CN}=\overrightarrow{NA}$.\\
			Theo quy tắc ba điểm ta có $\overrightarrow{BM}+\overrightarrow{CN}+\overrightarrow{AP}=\left(\overrightarrow{PN}+\overrightarrow{NA}\right)+\overrightarrow{AP}$  $=\overrightarrow{PA}+\overrightarrow{AP}=\overrightarrow0$.
			\item Theo  quy tắc ba điểm ta có
			\begin{eqnarray*}
				 \overrightarrow{OA}+\overrightarrow{OB}+\overrightarrow{OC}&=&\left(\overrightarrow{OP}+\overrightarrow{PA}\right)+\left(\overrightarrow{OM}+\overrightarrow{MB}\right)+\left(\overrightarrow{ON}+\overrightarrow{NC}\right)\\
				&=&\left(\overrightarrow{OM}+\overrightarrow{ON}+\overrightarrow{OP}\right)+\overrightarrow{PA}+\overrightarrow{MB}+\overrightarrow{NC}\\ 
				&=&\left(\overrightarrow{OM}+\overrightarrow{ON}+\overrightarrow{OP}\right)-\left(\overrightarrow{BM}+\overrightarrow{CN}+\overrightarrow{AP}\right). 
			\end{eqnarray*}
		\end{enumerate}
		Theo câu a)  \[\overrightarrow{BM}+\overrightarrow{CN}+\overrightarrow{AP}=\overrightarrow0.\] ta suy ra \[\overrightarrow{OA}+\overrightarrow{OB}+\overrightarrow{OC}=\overrightarrow{OM}+\overrightarrow{ON}+\overrightarrow{OP}.\]
	} 
\end{bt}
